\chapter*{Danh mục từ viết tắt và thuật ngữ}
\label{Danh mục từ viết tắt và thuật ngữ}

\begin{center}
    \begin{longtable}{|p{1cm}|p{2.75cm}|p{2.75cm}|p{8cm}|}
        \hline
        \textbf{Số thứ tự}  & \textbf{Từ viết tắt} & \textbf{Từ đầy đủ} & \textbf{Giải thích} \\
        \endfirsthead
        \hline
        \textbf{Số thứ tự}  & \textbf{Từ viết tắt} & \textbf{Từ đầy đủ} & \textbf{Giải thích} \\
        \endhead
        \hline
        01 & repository & repository & Là một nơi lưu trữ và quản lý mã nguồn, tài liệu, dữ liệu, và các tài nguyên khác của một dự án phần mềm. \\ 
        \hline
        02 & framework & framework & Là một kiến trúc hoặc nền tảng được dùng để xây dựng và phát triển các dự án phần mềm, bao gồm các công cụ, thư viện, và các quy tắc riêng. \\ 
        \hline
        03 & native & native & Là việc phát triển ứng dụng bằng các ngôn ngữ lập trình và công cụ được hỗ trợ chính thức bởi hệ điều hành mà ứng dụng đó nhắm tới, trong ngữ cảnh của báo cáo này là Android và iOS. \\ 
        \hline
        04 & SDK & Software Development Kit & Là một bộ công cụ phần mềm, cung cấp một tập hợp các công cụ, thư viện, và tài liệu dùng để phát triển ứng dụng cho một nền tảng cụ thể. \\ 
        \hline
        05 & package & package & Có thể hiểu như một thư viện, nó là tập hợp các mã nguồn, tài nguyên để phục vụ, hỗ trợ cho việc phát triển tính năng nào đó. \\ 
        \hline
        06 & API & Application Programming Interface & Là một tập hợp các quy tắc và giao thức cho phép các phần mềm khác tương tác và trao đổi thông tin với nhau. \\ 
        \hline
        07 & URL & Uniform Resource Locator & Là một chuỗi ký tự đại diện cho địa chỉ duy nhất của một tài nguyên trên mạng internet, thường dùng để truy cập trang web và các dịch vụ trực tuyến. \\ 
        \hline
        08 & SPA & Single Page Application & Là một kiểu ứng dụng web chỉ có một trang duy nhất. Các nội dung của trang được hiển thị và cập nhật mà không cần tải lại toàn bộ trang web. \\ 
        \hline
        09 & business logic & business logic & Là tập hợp các quy tắc và quy trình hoạt động của một tổ chức, hệ thống, hoặc ứng dụng nào đó, đại diện cho cách các đối tượng trên hoạt động và đưa ra quyết định. \\ 
        \hline
        10 & collection & collection & Là khái niệm tương đương với bảng trong cơ sở dữ liệu quan hệ. Chúng là tập hợp các thực thể có liên quan đến nhau, sở hữu một hoặc nhiều điểm chung nhất định nào đó. \\ 
        \hline
        11 & document & document & Là một khái niệm tương đương với các dòng trong một bảng của cơ sở dữ liệu quan hệ, thường lưu trữ thông tin bằng các cặp key-value. Các document trong cùng collection không nhất thiết phải cùng cấu trúc dữ liệu. \\ 
        \hline
        12 & field & field & Là một khái niệm tương đương với các cột (hay thuộc tính) trong một bảng của cơ sở dữ liệu quan hệ. Mỗi document gồm nhiều field, mỗi field có một tên định danh duy nhất và chứa một giá trị tương ứng. \\ 
        \hline
        13 & map & map & Là một kiểu dữ liệu được sử dụng để lưu trữ thông tin dưới dạng các cặp key-value. Mỗi key là duy nhất và chứa một giá trị nào đó. \\ 
        \hline
        14 & SEO & Search Engine Optimization & Là quá trình tối ưu hóa một trang web nhằm nâng cao khả năng xuất hiện trên các công cụ tìm kiếm, tăng khả năng tìm thấy và lượng truy cập vào trang web. \\ 
        \hline
        15 & prototype & prototype & Là phiên bản ban đầu của sản phẩm, cung cấp cái nhìn tổng quan, dùng để thử nghiệm, kiểm tra, và thu thập ý kiến trước khi triển khai chính thức. \\ 
        \hline
        16 & MVP & Minimum Viable Product & Là phiên bản đơn giản và tối thiểu của sản phẩm, với đủ các tính năng để hoạt động, nhằm giới thiệu và thu thập phản hồi từ người dùng. \\ 
        \hline
        17 & codebase & codebase & Là tập hợp mã nguồn, cách tổ chức thư mục và tệp tin, cách thức cấu hình, các thư viện có sẵn, và bất kỳ tài nguyên nào liên quan đến dự án. \\ 
        \hline
        18 & commit & commit & Là một thao tác quản lý mã nguồn bên trong Git, mô tả việc thực hiện gửi và xác nhận các thay đổi của mã nguồn vào nơi lưu trữ. \\ 
        \hline
        19 & rebase & rebase & Là một thao tác quản lý mã nguồn bên trong Git. Nó cho phép kết hợp và áp dụng các thay đổi từ một nhánh vào nhánh khác. \\ 
        \hline
        20 & pull request & pull request & Là một cách cho các thành viên trong nhóm xem xét và thảo luận về các thay đổi trong mã nguồn, trước khi hợp nhất chúng vào nhánh chính. \\ 
        \hline
        21 & scrum & scrum & Là một quy trình phát triển phần mềm theo mô hình Agile, sử dụng rộng rãi trong các dự án linh hoạt và yêu cầu sự thay đổi liên tục. \\ 
        \hline
        22 & sprint & sprint & Là một chu kỳ thời gian theo quy trình phát triển phần mềm scrum. Ở mỗi sprint, nhóm sẽ tập trung vào việc hoàn thành một tập hợp các mục tiêu, công việc cụ thể. \\ 
        \hline
        23 & story point & story point & Là một đơn vị đo lường được sử dụng để ước tính khối lượng, chi phí, độ phức tạp của một công việc, một nhiệm vụ cụ thể. \\ 
        \hline
        24 & QC & Quality Control & Là quá trình đảm bảo rằng sản phẩm, dịch vụ, hoặc công việc được giao đáp ứng các yêu cầu đã đặt ra từ trước. \\ 
        \hline
        25 & interface & interface & Là một tập hợp các phương thức (method) mà một lớp cần tự cài đặt (implement) nếu muốn sử dụng các phương thức của interface đó. \\ 
        \hline
        26 & CI/CD & Continuous Integration, Continuous Deployment & Việc tự động hóa trong quy trình phát triển sản phẩm và sau đó phân phối đến người dùng một cách tự động. \\ 
        \hline
        27 & container deployment & container deployment & Là quá trình đóng gói sản phẩm trước khi triển khai lên môi trường, hệ điều hành mong muốn. \\ 
        \hline
        28 & SSR & Server-Side Rendering & Là việc máy chủ sẽ truy vấn dữ liệu sao đó tích hợp vào HTML trước khi phản hồi về cho trình duyệt. \\ 
        \hline
        29 & CSR & Client-Side Rendering & Là việc máy chủ chỉ trả về các đoạn HTML, sau đó trình duyệt sẽ tiếp tục truy vấn dữ liệu thông qua việc gọi API và hiển thị lên cho người dùng. \\ 
        \hline
        30 & serverless & serverless & Là một mô hình phát triển và triển khai ứng dụng sử dụng điện toán đám mây để quản lý và vận hành, thay cho máy chủ truyền thống. \\
        \hline
        31 & extension & extension & Là một phần mở rộng của phần mềm, ứng dụng, hay dịch vụ nào đó, giúp bổ sung và mở rộng thêm tính năng.  \\ 
        \hline
        32 & SSL & Secure Sockets Layer & Là một giao thức bảo mật sử dụng để bảo vệ thông tin được truyền giữa máy chủ và trình duyệt hoặc giữa hai máy chủ. \\ 
        \hline
        33 & TLS & Transport Layer Security & Là một phiên bản nâng cấp của SSL. \\ 
        \hline
        34 & Cache & Cache & Là bộ nhớ đệm có chức năng lưu trữ tạm thời các nội dung thường được truy cập và ít thay đổi trong một khoảng thời gian cụ thể. \\ 
        \hline
        35 & HTTPS & Hypertext Transfer Protocol Secure & Là một giao thức truyền tải dữ liệu an toàn và bảo mật trên internet, và là phiên bản bảo mật của giao thức HTTP thông thường. \\ 
        \hline
        36 & APIs Server & APIs Server & Là một thành phần được thiết kế để xử lý yêu cầu từ các ứng dụng thông qua các API.  \\ 
        \hline
        37 & instance & instance & Trong AWS Elastic Compute Cloud, thuật ngữ "instance" được sử dụng để chỉ một máy ảo đang chạy trên nền tảng AWS. \\ 
        \hline
        38 & WebSocket & WebSocket & Là một giao thức truyền tải dữ liệu hai chiều (full-duplex) trên một kết nối mạng duy trì liên tục giữa máy khách (client) và máy chủ (server). \\
        \hline
        39 & module & module & Là một phần mã nguồn hoặc chức năng độc lập và có thể tái sử dụng. \\
        \hline
        \caption{Những từ viết tắt và thuật ngữ}
    \end{longtable}
\end{center}